\section{Publish a descriptor after the complete walk}
The parser's output is thirty-two bytes in the simulator's little-endian memory.
It records both the audio description and the location of its bytes. The following
layout is an interface contract; it is not a copy of a file header.

\tagpdfsetup{table/header-rows={1}}
\begin{center}
\begin{tabular}{rll}
Offset & Field & Width \\
\hline
0 & Frame rate & Unsigned 32 bits \\
4 & Channel count & Unsigned 32 bits \\
8 & Frame count & Unsigned 64 bits \\
16 & PCM offset from file start & Unsigned 64 bits \\
24 & PCM byte count & Unsigned 64 bits \\
\end{tabular}
\end{center}

For the worked file, the five values are 16000, two, four, 56 and sixteen. They refer
to two different extents: four frames in the sixteen-byte PCM slice, within an
eighty-two-byte file. Derive a PCM address by adding the stored offset to the particular
file buffer's full address. The offset does not retain an earlier buffer's address.

Retain candidate fields privately while scanning. Publish all thirty-two bytes only
after every chunk, including any chunk after data, has passed. On failure the caller's
descriptor stays unchanged. The bound, filter and encoder consume only the first
sixteen bytes, containing rate, channels and frames. They derive their own byte count
and never follow the trailing offset or length fields from an unrelated parsed file.
